\documentclass{article}
\usepackage{base}
\usepackage{url}

\title{Manfishing.net}
\author{Kyler Cain}
\date{2026-07-28}
% tags: unix, games, learning

\begin{document}
\maketitle

\begin{abstract}
Manfishing is a daily guessing game played on the Unix manual. Each round shows
the \verb|SEE ALSO| cross-references of a hidden man page and and you have to guess
which one. Ten rounds a day, the same ten for everyone, and includes the full text
of the manuals so you can read them after you play. Play it at at \url{https://manfishing.net/}.
\end{abstract}

\section{How it plays}

A round gives you a list of cross-references and nothing else, for example:

\begin{quote}
\verb|pldd(1)   sprof(1)   ld.so(8)   ldconfig(8)|
\end{quote}

Those four pages are what one particular manual page lists under \verb|SEE ALSO|.
Then you get to guess the name of the page that includes those options. 
Here the answer is \verb|ldd|\cite{ldd}, which you may be able to guess as the clues
are all about the dynamic linker.

There are 10 rounds each day, everyone shares them so you can compare your guesses
with others. When you finish a round, the reveal shows the page's one-line description
and a link to the full page, so you can read it and learn what you missed.

At the end of the game, you get a summary of how many you got right and wrong and an 
easily shareable emoji recap of your performance, similar to the one Wordle\cite{wordle}
and Catfishing.net\cite{catfishing} offer. 

\section{Why \texttt{SEE ALSO}}

\verb|SEE ALSO| is a conventional section in the man pages\cite{man_pages_7}, 
and is known to be human-generated and should be based on the author's expectation
of what other information would be most valuable to a reader. Even better is that,
thanks to the man-pages project\cite{man_pages_project}, there are thousands.

\section{What it is for}

My goal in making this game was to get people to read the manual, and to do it in a
way that is more fun than just reading and hopefully would involve a social component.
My coworkers and I have been playing the catfishing game\cite{catfishing} for a while,
and I wanted to make something similar for the Unix manual.

The game is deliberately structured to encourage further reading:

\begin{itemize}
  \item \textbf{Every round ends on a page.} Win or lose, the reveal links
  provides you a link to the man page for the potential answers.

  \item \textbf{Wrong-but-reasonable is not punished.} Many groups of clues fit
  several pages equally well. Nothing in the clues for \verb|latex(1)|
  distinguishes it from \verb|lualatex(1)| or \verb|xelatex(1)|, so all of them
  count, and the reveal names the whole family.

  \item \textbf{Randomized rounds.} The ten pages for each day are drawn at random from
  the pool of eligible pages, ensuring a fresh challenge every day, and there are so
  many pages that you're basically guaranteed to learn something new every day.
\end{itemize}

\section{The scoreboard and difficulty ratings}

The game keeps track of how well the community does on each page so you can compare
yourself against the global average. The pages also include a difficulty rating that you
can contribute to that I plan to use to make the game more interesting in the future, 
for example by ensuring a certain mix of easy and hard pages each day.

\section{Try it}

It is live at \url{https://manfishing.net/}\cite{manfishing}, and there is a new
edition every day.

\begin{thebibliography}{99}

\bibitem{manfishing}
Kyler Cain, ``manfishing,''
\url{https://manfishing.net/}

\bibitem{catfishing}
Catfishing, ``catfishing,''
\url{https://catfishing.net/}

\bibitem{wordle}
Josh Wardle, ``Wordle,''
\url{https://www.nytimes.com/games/wordle/}

\bibitem{man_pages_7}
``man-pages(7), conventions for writing Linux man pages,''
\url{https://man7.org/linux/man-pages/man7/man-pages.7.html}

\bibitem{man_pages_project}
Michael Kerrisk, ``The Linux man-pages project,''
\url{https://www.kernel.org/doc/man-pages/}

\bibitem{ldd}
``ldd(1), print shared object dependencies,''
\url{https://man7.org/linux/man-pages/man1/ldd.1.html}

\end{thebibliography}

\end{document}
